\documentclass[a4paper,11pt]{ctexart}
\usepackage[left=1.9cm,right=1.9cm,top=2.0cm,bottom=2.0cm]{geometry}
\usepackage{amsmath,amssymb}
\usepackage{enumitem}
\usepackage{fancyhdr}
\usepackage{lastpage}
\usepackage{listings}
\usepackage{xcolor}
\usepackage{array}

\pagestyle{fancy}
\fancyhf{}
\renewcommand{\headrulewidth}{0pt}
\fancyfoot[C]{第~\thepage~页\quad 共~\pageref{LastPage}~页}

\lstset{
  language=C,
  basicstyle=\ttfamily\small,
  columns=fullflexible,
  keepspaces=true,
  breaklines=true,
  showstringspaces=false,
  tabsize=4,
  aboveskip=3pt,
  belowskip=3pt,
  xleftmargin=1.5em,
  escapeinside={(*@}{@*)},
  commentstyle=\color{black!70}
}

\newcommand{\blankline}[1]{\underline{\hspace{#1}}}
\newcommand{\ansbox}{\hfill(\qquad\quad)}
\newcommand{\secbox}[2]{%
  \noindent\begin{tabular}{|>{\centering\arraybackslash}m{1.5cm}|>{\centering\arraybackslash}m{1.5cm}|}
    \hline
    分\quad 数 & \\ \hline
    评卷人 & \\ \hline
  \end{tabular}\hspace{0.8em}%
  \begin{minipage}[c]{\dimexpr\textwidth-4.2cm\relax}
    \textbf{#1}#2
  \end{minipage}\par\vspace{0.5em}%
}

\begin{document}

\begin{center}
  {\large\textbf{华中科技大学人工智能与自动化学院 2023～2024 学年第一学期}}\\[4pt]
  {\Large\textbf{“C 语言程序设计”考试试卷（A 卷）}}\\[8pt]
  考试方式：\underline{\quad 闭卷 \quad}\qquad
  考试日期：\underline{\quad 2024-01-18 \quad}\qquad
  考试时长：\underline{\quad 150 \quad} 分钟
\end{center}

\vspace{4pt}
\begin{center}
  \begin{tabular}{|>{\centering\arraybackslash}p{1.6cm}|*{7}{>{\centering\arraybackslash}p{1.5cm}|}}
    \hline
    题号 & 一 & 二 & 三 & 四 & 五 & 六 & 总分 \\ \hline
    分数 & & & & & & & \\[8pt] \hline
  \end{tabular}
\end{center}

\begin{center}
  \small 院系：\underline{\hspace{2.3cm}}\quad 班级：\underline{\hspace{2.0cm}}\quad 姓名：\underline{\hspace{2.0cm}}\quad 学号：\underline{\hspace{2.3cm}}\quad 成绩：\underline{\hspace{1.8cm}}
\end{center}
\vspace{6pt}

% ==================== 一 ====================
\secbox{一、判断下列语句或程序段的对错。}{（“$\times$”表示错，“$\checkmark$”表示对）（10 分）}

\begin{enumerate}[label=(\arabic*), leftmargin=2.4em, itemsep=3pt]
  \item \texttt{char x=0, y=x ;} \ansbox
  \item \texttt{int  xyz-2 = 1;} \ansbox
  \item \texttt{printf("\%d\textbackslash n",(40,50,60));} \ansbox
  \item \texttt{float x[10];}\\
        \texttt{int *p = x;} \ansbox
  \item \texttt{register int x;}\\
        \texttt{scanf("\%d",\&x);} \ansbox
  \item \texttt{char  str[20];}\\
        \texttt{str="hust aia";} \ansbox
  \item \texttt{int  *p , a[2][3];}\\
        \texttt{p = *a+1;} \ansbox
  \item \texttt{int  s[]= \{5,6,7,8\} ;} \ansbox
  \item \texttt{char  **pp, a[4][5];}\\
        \texttt{pp = a;} \ansbox
  \item \texttt{int x='\textbackslash 107' ;} \ansbox
\end{enumerate}

\vspace{8pt}

% ==================== 二 ====================
\secbox{二、计算下列表达式的值}{（10 分）}

\noindent 假设 \texttt{int} 和 \texttt{unsigned int} 类型均为 16 位长度，且各题彼此独立\\
设 \texttt{unsigned int  a=10, b=17,c=5,d=3;}\\
\phantom{设 }\texttt{float f ;}

\begin{enumerate}[label=(\arabic*), leftmargin=2.4em, itemsep=3pt]
  \item \texttt{f = b/c;} \ansbox
  \item \texttt{!(a+b)+c\&\&b+c/2;} \ansbox
  \item \texttt{(a\^{}b) + ( c >>1|d)} \ansbox
  \item \texttt{f=(a-b>0?c/d:c\%d);} \ansbox
  \item \texttt{a=2,b=a* ++b;} \ansbox
\end{enumerate}

\vspace{8pt}

% ==================== 三 ====================
\secbox{三、程序改错（10 分）}{要求：不得改变程序框架，不得重写程序，无需文字说明，直接在代码上添加、删除和修改。}

\noindent 1、求 10 个浮点数的最大值和平均值，其中在 \texttt{main()} 函数中完成输入和显示求得的结果，在 \texttt{fun} 函数中求最大值和平均值。
\begin{lstlisting}
#include<stdio.h>

void  main()
{
    int i;
    float  x[10], max, average;

    for(i=0;i<10;i++)
        scanf("%f",x+i);
    average=fun(x,10,max);
    pirintf ("max=%f,average=%f\n", max, average) ;
}
fun(float *p, int n, int *pmax)
{
    float temp;
    int i;

    for(i=0;i<n;i++)
    {
        if (*pmax<p[i])
            *pmax=p[i];
        temp+=p[i];
    }
    return temp/n;
}
\end{lstlisting}

\noindent 2、求 5 个字符串中的最小字符串。
\begin{lstlisting}
#include <stdio.h>

void main()
{
    char *pstr[5]={"pascal","basic","prolog","cobol","lisp"};
    char *pmin;
    pmin=getminstr(*pstr,5);
    printf("min string=%s\n",pmin)
}
char getminstr(char **str, int n)
{
    char *temp,i;

    for(i=0;i<n;i++)
        if(temp>str[i])
            temp=str[i];

    return temp;
}
\end{lstlisting}

\vspace{8pt}

% ==================== 四 ====================
\secbox{四、程序填空}{（10 分）}

\noindent (1) 求 10--1000 某个闭区间的所有素数并在屏幕上显示，区间由用户从键盘输入。
\begin{lstlisting}
#include<stdio.h>

void main()
{
    int i,j;
    int a, b ;

    printf("please input a ,b:");
    scanf("%d%d", (*@\blankline{4.5cm}@*) );

    for ( (*@\blankline{4.5cm}@*) ;i++)
    {
        for (j=2; j<i ; j++)
        {
            if( (*@\blankline{4.5cm}@*) )
                (*@\blankline{3.5cm}@*) ;
        }
        if( (*@\blankline{3.5cm}@*) )
            printf( "%d\n" ,i);
    }
}
\end{lstlisting}

\noindent (2) 将某一个整数转换成对应的字符串.
\begin{lstlisting}
#include <stdio.h>
(*@\blankline{5.5cm}@*)
void main()
{
    int  n,  sign ;
    char str[10], i, j;

    scanf( "%d" ,&n);
    if((sign=n)<0)
        (*@\blankline{5.0cm}@*)
    i=0;
    do
    {
        (*@\blankline{5.0cm}@*)
    }(*@\blankline{4.5cm}@*)
    if(sign<0)
        str[i++]='-';
    (*@\blankline{5.0cm}@*)
    for(i=0,j=strlen(str)-1;i<j;i++,j--)
        n=str[i],str[i]=str[j],str[j]=n;
    puts(str);
}
\end{lstlisting}

\vspace{8pt}

% ==================== 五 ====================
\secbox{五、输出程序运行结果}{（25 分，每题 5 分，结果写在题目的右边）}

\noindent 1、
\begin{lstlisting}
#include <stdio.h>
void fun(int *, int);
void main()
{
    int a[5]={1, 2, 3, 4, 5}, i;

    fun(a+1, 3);
    for(i=1;i<4;i++)
        printf("%d\n",a[i]);
}
void fun(int *b,int n)
{
    int i;
    for(i=0;i<n;i++)
        b[i]=b[i]*b[i+1];
    return ;
}
\end{lstlisting}

\vspace{4pt}
\noindent 2、
\begin{lstlisting}
#include <stdio.h>
void main()
{
    int i, j, temp;
    int row=0, column=0;
    int a[3][3]={{6, 3, 1},{-4, 3, 2},{-9, 4, 7}};

    temp=a[0][0];
    for(i=0;i<3;i++)
        for(j=0; j<3; j++)
        {
            if(a[i][j]>temp)
            {
                temp=a[i][j];
                row=i+1;
                column=j+1;
            }
        }
    printf("temp=%d, row=%d, column=%d\n", temp, row, column);
}
\end{lstlisting}

\vspace{4pt}
\noindent 3、
\begin{lstlisting}
#include <stdio.h>
int n;
void func();
void main()
{
    static int x=3;
    int y;

    y=n+1;
    printf("main:x=%d, y=%d, n=%d\n", x, y, n);
    func();
    printf("main:x=%d, y=%d, n=%d\n", x, y, n);
}
void func()
{
    static int x;
    int y=10;
    x=x+2;
    n=n+2;
    y=y-n;
    printf("func:x=%d, y=%d, n=%d\n", x, y, n);
}
\end{lstlisting}

\vspace{4pt}
\noindent 4、
\begin{lstlisting}
#include <stdio.h>
#include <string.h>
struct person
{
    char name[20];
    int count;
};
void main()
{
    struct person leader[3]={{"li",5},{"zhang",6},{"wang",7}};
    char name[20],m,n;
    for(m=0;m<6;m++)
    {
        gets(name);
        for(n=0;n<3;n++)
            if(!strcmp(name,leader[n].name))
            {
                leader[n].count++;
                break;
            }
    }
    printf("\n");
    for(m=0;m<3;m++)
        printf("%s:%d\n",leader[m].name,leader[m].count);
}
\end{lstlisting}
\noindent 当程序运行时，键入情况如下：\\
\texttt{li}\\
\texttt{wang}\\
\texttt{lei}\\
\texttt{li}\\
\texttt{wang}\\
\texttt{zhang}\\
写出程序的输出结果。

\vspace{6pt}
\noindent 5、
\begin{lstlisting}
#include <stdio.h>
#include <string.h>
void main()
{
    char str[4][20]={ "capital", "indexing", "larger", "smaller"};
    int a,b;
    char temp[20];

    for(a=0;a<3;a++)
        for(b=0;b<3-a;b++)
            if(strcmp(str[b],str[b+1])>0)
            {
                strcpy(temp,str[b]);
                strcpy(str[b],str[b+1]);
                strcpy(str[b+1],temp);
            }
    for(a=1;a<3;a++)
        printf("%s\n",str[a]+a);
}
\end{lstlisting}

\vspace{8pt}

% ==================== 六 ====================
\secbox{六、编程（35 分）}{注意：不得使用全局变量，注意程序结构}

\noindent 1、某门课程最终成绩分别 \texttt{'A'}、\texttt{'B'}、\texttt{'C'}、\texttt{'D'}，90 分以上（含 90 分）为 \texttt{'A'}，70-89 分为 \texttt{'B'}，60-69 分为 \texttt{'C'}，60 分以下为 \texttt{'D'}。现要求从键盘输入成绩的等级，在屏幕上输出该成绩等级对应的成绩范围（8 分）

\vspace{6pt}
\noindent 2、编写一个函数求一个二维数组列累加和的最大值并返回列号；另编写 \texttt{main} 测试函数完成二维数组数的输入，函数调用，输出列累加和的最大值以及列号。（9 分）

\vspace{6pt}
\noindent 3、输入一行包含若干单词（假定不超过 10 个）的字符串，单词之间用空格分开，要求删除该行字符串中最大的那个单词（假定最大单词只有一个）形成新的字符串，并显示。（不考虑单词的合法性以及语句的语法）（9 分）

\vspace{6pt}
\noindent 4、编程处理某班 30 个学生 4 门课的成绩，它们分别是数学、物理、英语和计算机，按编号从小到大的顺序依次输入学生的姓名、性别和四门课的成绩，计算每个学生的平均成绩，并以清晰的格式从低分到高分的顺序打印平均分低于 80 分的学生成绩（9 分）

\end{document}
